\documentclass[12pt]{article}
\makeatletter
\def\input@path{{../../}}
\makeatother
\usepackage{sbc-template}
\makeatletter
\renewcommand{\inst}[1]{}
\makeatother
\usepackage{graphicx,url}
\usepackage[utf8]{inputenc}
\usepackage[T1]{fontenc}
\usepackage[brazil]{babel}
\usepackage{longtable,booktabs,array}
\providecommand{\tightlist}{%%
  \setlength{\itemsep}{0pt}\setlength{\parskip}{0pt}}
\sloppy
\title{Fundamentos de Reservoir Computing: da intuicao ao modelo minimo}
\author{}
\address{}
\begin{document}

\maketitle

\begin{resumo}
Reservoir Computing (RC) e um paradigma de modelagem temporal baseado em uma ideia simples: em vez de treinar toda a dinamica recorrente, mantemos um sistema dinamico rico e fixo, e treinamos apenas uma camada de leitura. Essa estrategia reduz a dificuldade de treino, preserva memoria temporal e produz uma representacao expandida da entrada ao longo do tempo. Neste artigo, apresentamos RC com o minimo de formalismo necessario para um leitor que vem do problema de previsao de demanda introduzido no artigo anterior. Mostramos o papel de entrada, estado interno, nao linearidade, washout e readout, e conectamos cada componente a um pipeline concreto de forecasting no Favorita. Ao final, o leitor entende por que RC pode ser util em series temporais e esta preparado para rodar um primeiro experimento reproduzivel.
\end{resumo}

\section{1. Introducao}\label{1-introducao}

Quando alguem aprende series temporais, uma das primeiras intuicoes e que o passado importa. Quando aprende redes neurais recorrentes, descobre que incorporar passado em um modelo pode ser poderoso, mas tambem caro, instavel e dificil de treinar. Reservoir Computing aparece como uma resposta elegante a esse impasse.

A ideia central de RC nao e "treinar menos por conveniencia". A ideia e explorar um fato estrutural: em muitos problemas temporais, boa parte do valor do modelo esta em construir uma representacao dinamica rica da entrada. Se essa representacao for suficientemente informativa, um readout linear simples pode ser suficiente para produzir boas previsoes.

Essa e a intuicao que vamos desenvolver neste artigo.

\section{2. A intuicao de um reservatorio}\label{2-a-intuicao-de-um-reservatorio}

Pense no reservatorio como um sistema dinamico que recebe a entrada ao longo do tempo e reage a ela de forma distribuida. Cada nova observacao nao e tratada isoladamente. Ela altera o estado interno do sistema, e esse estado carrega tracos do passado recente.

Em termos didaticos, o reservatorio faz tres coisas ao mesmo tempo:

\begin{itemize}
\tightlist
\item
  mistura a entrada atual com a memoria do passado;
\item
  aplica uma transformacao nao linear;
\item
  projeta a serie observada em um espaco de features mais rico.
\end{itemize}

O papel do readout e muito mais simples: receber esse estado expandido e aprender um mapeamento para a saida desejada.

Essa separacao de responsabilidades e a marca registrada de RC:

\begin{itemize}
\tightlist
\item
  a dinamica do reservatorio e fixa;
\item
  o readout e a parte treinada.
\end{itemize}

\section{3. O modelo minimo}\label{3-o-modelo-minimo}

Uma formulacao minima e suficiente para a serie.

Se \texttt{u\_t} e a entrada no instante \texttt{t}, \texttt{x\_t} e o estado do reservatorio, e \texttt{y\_t} e a previsao, podemos escrever:

\texttt{x\_t\ =\ (1\ -\ a)\ x\_\{t-1\}\ +\ a\ *\ f(W\_in\ u\_t\ +\ W\ x\_\{t-1\}\ +\ b)}

\texttt{y\_t\ =\ W\_out\ {[}1;\ u\_t;\ x\_t{]}}

onde:

\begin{itemize}
\tightlist
\item
  \texttt{W\_in} projeta a entrada no reservatorio;
\item
  \texttt{W} define a recorrencia interna;
\item
  \texttt{b} e um termo de vies;
\item
  \texttt{f} e uma nao linearidade, tipicamente \texttt{tanh};
\item
  \texttt{a} e o leak rate, que controla o quanto o estado novo substitui o estado antigo;
\item
  \texttt{W\_out} e o readout treinado.
\end{itemize}

No projeto desta serie, o readout e uma regressao Ridge. Isso nao e um detalhe menor. E uma escolha pedagogica importante: manter a leitura simples ajuda a mostrar que o ganho do metodo vem da representacao dinamica, nao de uma cabeca preditiva sofisticada.

\section{4. O que e treinado e o que nao e treinado}\label{4-o-que-e-treinado-e-o-que-nao-e-treinado}

Essa e a pergunta que mais ajuda iniciantes a organizar a intuicao.

Em RC classico:

\begin{itemize}
\tightlist
\item
  \texttt{W\_in}, \texttt{W} e \texttt{b} normalmente sao inicializados e mantidos fixos;
\item
  apenas \texttt{W\_out} e ajustado aos dados.
\end{itemize}

Isso produz duas consequencias didaticas:

\begin{enumerate}
\def\labelenumi{\arabic{enumi}.}
\tightlist
\item
  o treino tende a ser mais estavel e barato do que treinar toda uma RNN;
\item
  a qualidade do reservatorio passa a depender fortemente de como sua dinamica foi desenhada.
\end{enumerate}

Ou seja: RC troca parte da dificuldade de otimizacao por uma dificuldade diferente, que e projetar e interpretar o reservatorio.

\section{5. Como o tempo entra no modelo}\label{5-como-o-tempo-entra-no-modelo}

RC so faz sentido se entendermos memoria. Em forecasting, memoria significa que a previsao atual depende de informacao passada que ainda esta representada no estado do sistema.

No contexto desta serie, isso aparece de tres formas:

\begin{itemize}
\tightlist
\item
  a demanda de hoje depende da demanda recente;
\item
  promocao atual pode ter efeito condicionado pelo historico recente;
\item
  padroes de calendario afetam a previsao com periodicidade curta e longa.
\end{itemize}

O estado \texttt{x\_t} serve justamente para condensar esse historico. Ele nao guarda uma copia literal do passado. Ele guarda uma transformacao dinamica do passado.

\subsection{5.1 Washout}\label{51-washout}

No inicio da serie temporal, o estado do reservatorio ainda esta fortemente influenciado pela inicializacao. Por isso, e comum descartar os primeiros passos antes de treinar o readout. Esse periodo e chamado de \texttt{washout}.

No codigo do projeto, o RC atual usa \texttt{washout\ =\ 14}. Essa escolha tem interpretacao simples: permitimos que o reservatorio "esqueca" a inicializacao antes de confiar em seu estado como feature preditiva.

\section{6. Um pseudocodigo minimo}\label{6-um-pseudocodigo-minimo}

\begin{verbatim}
inicializar W_in, W, b
inicializar estado x = 0

para cada tempo t no treino:
    receber entrada u_t
    atualizar estado x_t
    se t > washout:
        armazenar feature [1, u_t, x_t]
        armazenar alvo y_t

treinar readout linear sobre features e alvos

para cada tempo t no teste:
    receber entrada u_t
    atualizar estado x_t
    prever y_t com o readout
\end{verbatim}

Esse pseudocodigo explica por que RC e tao atraente pedagogicamente. O pipeline e simples o suficiente para ser implementado do zero e rico o suficiente para introduzir ideias profundas sobre dinamica temporal.

\section{7. Como o projeto implementa RC no caso Favorita}\label{7-como-o-projeto-implementa-rc-no-caso-favorita}

No repositorio desta serie, o RC classico esta implementado em \texttt{code/rc/model.py}.

O vetor de entrada sequencial e construindo a partir de sete componentes:

\begin{itemize}
\tightlist
\item
  demanda previa normalizada;
\item
  promocao normalizada;
\item
  \texttt{dow\_sin} e \texttt{dow\_cos};
\item
  \texttt{month\_sin} e \texttt{month\_cos};
\item
  indicador de fim de semana.
\end{itemize}

Essa escolha e pedagogicamente feliz por dois motivos:

\begin{enumerate}
\def\labelenumi{\arabic{enumi}.}
\tightlist
\item
  ela preserva interpretabilidade;
\item
  ela permite que o reservatorio trabalhe sobre um conjunto pequeno, mas informativo, de sinais.
\end{enumerate}

O reservatorio usado no projeto segue a forma minima discutida acima:

\begin{itemize}
\tightlist
\item
  \texttt{n\_reservoir\ =\ 80}
\item
  \texttt{spectral\_radius\ =\ 0.6}
\item
  \texttt{leak\_rate\ =\ 0.15}
\item
  \texttt{input\_scale\ =\ 0.8}
\item
  \texttt{bias\_scale\ =\ 0.1}
\item
  \texttt{washout\ =\ 14}
\item
  \texttt{ridge\_alpha\ =\ 1e-3}
\end{itemize}

Mesmo sem esgotar o espaco de hiperparametros, essa configuracao ja permite mostrar como um reservatorio opera em um problema realista.

\section{8. Por que RC pode ser util para previsao de demanda}\label{8-por-que-rc-pode-ser-util-para-previsao-de-demanda}

Em previsao de demanda, RC e interessante porque:

\begin{itemize}
\tightlist
\item
  lida naturalmente com dependencia temporal;
\item
  incorpora memoria sem exigir treino recorrente completo;
\item
  pode combinar sinais autoregressivos e exogenos;
\item
  oferece uma ponte conceitual clara para reservatorios fisicos e, depois, quanticos.
\end{itemize}

Isso nao quer dizer que RC seja automaticamente superior a modelos classicos. Pelo contrario: parte do valor didatico da serie e mostrar que RC precisa ser comparado de forma honesta com baselines fortes.

Mas, como paradigma, RC e poderoso porque ensina uma forma diferente de pensar o problema: antes de perguntar "qual e o modelo mais complexo?", perguntamos "qual dinamica produz uma representacao temporal util?".

\section{9. Limites e promessas}\label{9-limites-e-promessas}

RC nao resolve tudo. Um reservatorio ruim continua sendo ruim, mesmo com readout simples. Alem disso:

\begin{itemize}
\tightlist
\item
  hiperparametros importam;
\item
  memoria insuficiente prejudica o modelo;
\item
  memoria excessiva pode tornar a dinamica lenta ou confusa;
\item
  reservatorios pequenos podem subrepresentar o problema;
\item
  reservatorios grandes podem aumentar custo sem ganho real.
\end{itemize}

Essas limitacoes nao diminuem o valor do paradigma. Elas abrem as perguntas certas para os proximos artigos:

\begin{itemize}
\tightlist
\item
  como construir um primeiro pipeline reproduzivel?
\item
  como ler memoria e nao linearidade no comportamento do modelo?
\item
  como comparar RC com outros metodos em protocolo justo?
\end{itemize}

\section{10. Conclusao}\label{10-conclusao}

Neste artigo, Reservoir Computing foi apresentado como uma ideia antes de ser tratado como um algoritmo. O leitor viu que RC separa dinamica e leitura, usa um reservatorio fixo para produzir features temporais e treina apenas um readout simples.

Essa formulacao e suficiente para o proximo passo da serie: sair do formalismo minimo e construir um experimento real, reproduzivel, com o dataset Favorita.

\section{Entregaveis associados no repositorio}\label{entregaveis-associados-no-repositorio}

\begin{itemize}
\tightlist
\item
  implementacao do RC: \texttt{code/rc/model.py}
\item
  utilitarios sequenciais: \texttt{code/common/sequential\_features.py}
\item
  benchmark consolidado: \texttt{code/RESULTS.md}
\end{itemize}

\section{Referencias}\label{referencias}

\begin{enumerate}
\def\labelenumi{\arabic{enumi}.}
\tightlist
\item
  Jaeger, H. The "echo state" approach to analysing and training recurrent neural networks.
\item
  Maass, W., Natschlager, T., Markram, H. Real-time computing without stable states.
\item
  Lukosevicius, M., Jaeger, H. Reservoir computing approaches to recurrent neural network training.
\end{enumerate}

\section*{Resultados Computacionais Associados}

Os resultados computacionais associados a este artigo estao organizados na subpasta \texttt{computational\_results\_20260402\_222902/}, localizada dentro desta pasta \LaTeX. Esse diretorio contem o arquivo \texttt{REPORT.md}, tabelas em CSV e imagens comparativas apropriadas ao escopo do artigo.

\end{document}
